\documentclass[UTF8, fontset=fandol,12pt, a4paper]{ctexart}

\usepackage{amsmath, amssymb}
\usepackage{pgfplots}
\usepackage{tikz}
\pgfplotsset{compat=1.18}
\usepackage[margin=2.5cm]{geometry}

\title{信号与系统分析-hw1}
\author{黄泱博 \quad 2024012867}
\date{\today}

\begin{document}
\maketitle

\section* {第一题}

\subsection{$x_2(t) = e^{(-1+j)t}$}

$x_2(t)$ 不是周期信号，因为$\lvert x_2(t) \rvert$ = $e^{-t}$ 随时间单调递减。

\subsection{$x_3[n] = e^{j7\pi n}$}

$x_3[n]$ 是周期信号， 因为$\omega = 7\pi, \frac{\omega}{2\pi} = \frac{7}{2} \in \mathbb{Q}$，由上亦有：基波周期 $T_0 = 2$

\subsection{$x_5[n] = 3e^{j3/5(n+1/2)}$}

$x_5[n] = 3e^{j\frac{3}{5}n} \times e^{j\frac{3}{10}}$， 后者为const，只需要关注前者的周期性即可。

因为$\omega = \frac{3}{5}, \frac{2\pi}{\omega} = \frac{10\pi}{3} \notin \mathbb{Q}$， 因此有$x_5[n]$ 不为周期信号。

\section* {第二题}

$x_1(t) = 2\cos{(10t+1)}$ 的基波周期 $T_1 = \frac{2\pi}{10} = \frac{\pi}{5}$, $x_2(t) = -\sin{(4t-1)}$ 的基波周期 $T_2 = \frac{2\pi}{4} = \frac{\pi}{2}$。 因此， $x(t) = x_1(t) + x_2(t)$ 的基波周期为 $T_0 = k_1 T_1 = k_2 T_2 (k_1, k_2 \in \mathbb{N}) $ 于是有 $\frac{k_1}{k_2} = \frac{T_2}{T_1} = \frac{5}{2}$，因此有$T_0 = \pi$

\section* {第三题}

\subsection {输出信号}

\begin{align}
  y[n] &= y_1[n-2] + \frac{1}{2} y_1[n-3] \\
       &= (2x[n-2] + 4x[n-3]) \\
       &\quad + \frac{1}{2}\bigl(2x[n-3] + 4x[n-4]\bigr) \\
       &= 2x[n-2] + 5x[n-3] + 2x[n-4].
\end{align}

\subsection {交换$S_1$、$S_2$次序}

\begin{align}
  y'[n] &= 2y_2[n] + 4y_2[n-1] \\
        &= 2(x[n-2] + \frac{1}{2}x[n-3]) \\
        &\quad + 4(x[n-3] + \frac{1}{2}x[n-4]) \\
        &= 2x[n-2] + 5x[n-3] + 2x[n-4].
\end{align}

由上可得：交换次序后输入-输出关系不变。

\section* {第四题}

\subsection{是否有记忆性}

$y[n] = x[n]x[n-2]$ 依赖于过往的输入值 $x[n-2]$， 因而有记忆。

\subsection{输入为 $A\delta (n)$ 时的输出}

$$y[n] = (A\delta [n])(A\delta [n-2]) = A^2 \delta[n]\delta[n-2] = 0$$

\subsection{是否可逆}

由上，系统显然不可逆。

\section* {第五题}

\subsection{因果性}

系统不是因果的，因为当 $t < 0$ 时，$t < \sin{t}$， 而 $y(t) = x(\sin{t})$ 依赖未来值。

\subsection{线性}

系统是线性的，因为对于输入 $x(t) = ax_1(t) + bx_2(t)$：
$$y(t) = x(\sin{t}) = ax_1(\sin{t}) + bx_2(\sin{t}) = ay_1(t) + by_2(t), \forall a, b \in \mathbb{C}$$

\section* {第六题}

\subsection{$x(t-1)$}

\begin{tikzpicture}
  \begin{axis}[
    axis lines=middle,
    xlabel={$x$}, ylabel={$y$},
    xmin=-3, xmax=4,
    ymin=-2, ymax=3,
    xtick={-2,-1,0,1,2,3},
    ytick={-1,1,2},
    width=7cm, height=5cm,
    scale only axis,
    tick label style={font=\small},
  ]
    \addplot[
      blue, thick,
      mark=*, mark size=2pt,
      mark options={fill=red, draw=red}
    ] coordinates {
      (-1,0) (-1,-1) (0,0) (0,1) (1,1)
      (1,2) (2,2) (2,1) (3,0)
    };
  \end{axis}
\end{tikzpicture}

\subsection{$x(4-\frac{t}{2})$}

\begin{tikzpicture}
  \begin{axis}[
    axis lines=middle,
    xlabel = {$x$}, ylabel = {$y$},
    xmin = -2, xmax = 15,
    ymin = -2, ymax = 3.5,
    width = 7cm, height = 5cm,
    scale only axis,
    tick label style={font=\small},
  ]
    \addplot[
      blue, thick,
      mark=*, mark size=2pt,
      mark options={fill=red, draw=red}
    ] coordinates {
      (4,0) (6,1) (6,2) (8,2) (8,1) (10,1) (10,0) (12,-1) (12,0)
    };
  \end{axis}
\end{tikzpicture}

\subsection{$x(t)[\delta(t+\frac{3}{2})-\delta(t-\frac{3}{2})]$}

\begin{tikzpicture}
  \begin{axis}[
    axis lines=middle,
    xlabel={$x$}, ylabel={$y$},
    xmin=-2, xmax=2,
    ymin=-1, ymax=2,
    width=7cm, height=5cm,
    scale only axis,
    tick label style={font=\small},
  ]
    \addplot[
      only marks,
      mark=*, mark size=2pt,red,
    ] coordinates{
      (-3/2, -1/2) (3/2, -1/2)
    };

    \node[below left] at (axis cs:-1, -1/2) {$(-\frac{3}{2}, -\frac{1}{2})$};
    \node[below right] at (axis cs: 1, -1/2) {$(\frac{3}{2}, -\frac{1}{2})$};
    
  \end{axis}
\end{tikzpicture}

\section* {第七题}

\subsection{$x(t) = e^{j(\pi t - 1)}$}

$x(t) = e^{j \pi t} e^{-j}$ 为周期信号， 基波周期 $T_0 = \frac{2\pi}{\omega} = 2$

\subsection{$x(t) = [\cos{(2t - \frac{\pi}{3})}]^2$}

$x(t) = \frac{cos(4t - \frac{2\pi}{3}) + 1}{2}$ 为周期信号， 基波周期 $T_0 = \frac{2\pi}{\omega} = \frac{\pi}{2}$

\subsection{$x(t) = Ev(\sin{(4\pi t)}u(t))$}

\[
  x(t) = \frac{\sin{(4\pi t)u(t)} + \sin{(-4\pi t)}u(-t)}{2} = \begin{cases}
    \frac{\sin{(4\pi t)}}{2}, & t \ge 0, \\
    -\frac{\sin{(4\pi t)}}{2}, & t < 0.
  \end{cases}
\]

0 是 $x(t)$ 唯一的导数不连续点。 故不是周期函数。

\section* {第八题}

\subsection{$y(t) = x(t-2) + x(2-t)$}

\begin{enumerate}
  \item 不是无记忆的，因为不只依赖于当前时间的值。
  \item 是时不变的，因为 $\forall t_0 \in \mathbb{R}$, 有
        $$x_2(t) = x_1(t - t_0)$$
        $$y_2(t) = x_2(t - 2) + x_2(2 - t) = x_1(t - t_0 - 2) + x_1(2 + t_0 - t) = y_1(t - t_0)$$
  \item 是线性的， 因为$\forall a, b \in \mathbb{R}$, 有
        \begin{align}
          y(t) &= (ax_1+bx_2)(t-2) + (ax_1+bx_2)(2-t) \\
               &= a(x_1(t-2) + x_1(2-t)) + b(x_2(t-2)+x_2(2-t)) \\
               &= ay_1(t)+by_2(t)
        \end{align}
  \item 不是因果的，因为当$t=0$时，$y(0)=x(-2)+x(2)$依赖于未来值。
  \item 是稳定的，因为若$\lvert x(t) \rvert < M, \forall t \in \mathbb{R}$ 则有：
        $$\lvert y(t) \rvert < \lvert x(t-2) \rvert + \lvert x(2-t) \rvert < 2M$$
\end{enumerate}

\subsection{$y(t)=\int_{-\infty}^{2t}x(\tau)d\tau$}
\begin{enumerate}
  \item 显然有记忆，因为不只依赖于当前时间的值
  \item 不是时不变的，因为$\forall t_0 \in \mathbb{R}, t_0 \ne 0$
        $$x_2(t) = x_1(t - t_0)$$
        $$y_2(t) = \int_{-\infty}^{2t}x_2(\tau)d\tau = \int_{-\infty}^{2t}x_1(\tau - t_0)d\tau = \int_{-\infty}^{2t - t_0} x_1(\tau) d\tau \ne y_1(t-t_0) $$
  \item 是线性的，因为积分是线性算子。
  \item 不是因果的，对于$t>0$ 时会依赖于未来值。
  \item 不是稳定的， 取 
        \[
          x(t) = \begin{cases}
            0, & t \le 0, \\
            1, & t > 0.
          \end{cases}
        \]
        显然$x(t)$有界而$y(t)$无限增长。
\end{enumerate}

\subsection{
  \[
    y(t) = \begin{cases}
      0, &t < 0, \\
      x(t) + x(t - 2). &t \ge 0
    \end{cases}
  \]
}

\begin{enumerate}
  \item 显然有记忆，因为不只依赖于当前时间的值
  \item 不是时不变的，不妨取$x_1(t) \equiv 1, t_0 > 0$
        \[
          y_2(t) = \begin{cases}
            0, &t<0, \\
            2, &t\ge 0.
          \end{cases}
        \]
        而
        \[
          y_1(t - t_0) = \begin{cases}
            0, &t<t_0, \\
            2, &t\ge t_0.
          \end{cases}
        \]
        显然$y_1(t-t_0) \ne y_2(t)$
  \item 是线性的，因为$y(t)$是关于$t$分段线性的。
  \item 是因果的，因为只依赖于现在即过往输入。
  \item 是稳定的，证明类似于第一题的稳定性证明.
\end{enumerate}

\subsection{$y(t) = \frac{dx(t)}{dt}$}
\begin{enumerate}
  \item 有记忆的，不止取决于当前输入值
  \item 是时不变的，因为对于$\forall t_0 \in \mathbb{R}$
        $$x_2(t) = x_1(t-t_0)$$
        $$y_2(t) = \frac{dx_2(t)}{dt} = \frac{dx_1(t-t_0)}{dt} = y_1(t-t_0)$$
  \item 是线性的，因为微分是线性算子
  \item 取决于实现方式，若使用左导数，可以认为是因果的，若采用双侧导数存在且相等，则不是因果的。
  \item 不是稳定的，取$x(t) = \sin{(t^2)}$,则$y(t) = 2t \cos{(t^2)}$ 则有: \\
        $$\lim_{n \rightarrow \infty}2\sqrt{2n\pi}\cos{(2n\pi)} = \infty$$
\end{enumerate}

\section* {第九题}

\subsection{$y[n] = Re(x[n])$可加}

显然，因为$Re(x_1 + x_2) = Re((a_1+jb_1) + (a_2+jb_2)) = a_1+a_2 = Re(x_1) + Re(x_2), \forall x_1, x_2 \in \mathbb{C}$。

\subsection{$y[n] = Re(e^{\frac{j\pi n}{4}}x[n])$是否可加}

依然可加， 因为:
\begin{align}
  & \quad Re((u + jv)((a_1+jb_1)+(a_2+jb_2))) \\
  & = u(a_1+a_2)-v(b_1+b_2) \\
  & = Re((u+jv)(a_1+jb_1)) + Re((u+jv)(a_2+jb_2))
\end{align}

\section* {第十题}

已知线性时不变系统对$x_1(t)$的响应$y_1(t)$

\subsection{绘制系统对$x_2(t)$的响应}

\begin{center}
\begin{tikzpicture}[scale=1.5]
  \draw[->] (-1,0) -- (5,0) node[right] {$x$};
  \draw[->] (0,-3) -- (0,3) node[above] {$y$};

  \draw[thick, blue]
    (0,0) -- (1,2) -- (2,0) -- (3,-2) -- (4, 0);

  \foreach \x in {1,2,3,4}
    \draw (\x,0.05) -- (\x, -0.05)
      node[below] {$\x$};

  \node[below left] at (0,0) {$0$};
  
  \foreach \y in {-2,-1,1,2}
    \draw (-0.05,\y) -- (0.05,\y)
      node[left] {$\y$};
  
\end{tikzpicture}
\end{center}

\subsection{绘制系统对$x_3(t)$的响应}

\begin{center}
\begin{tikzpicture}[scale=1.5]
  \draw[->] (-2,0) -- (3,0) node[right] {$x$};
  \draw[->] (0,-1) -- (0,3) node[above] {$y$};

  \draw[thick,blue]
    (-1,0) -- (0,2) -- (1,2) -- (2,0);

  \foreach \x in {-1,1,2}
    \draw (\x,-0.05) -- (\x,0.05)
      node[below] {$\x$};

  \node[below, left] at (0,-0.15) {$0$};

  \foreach \y in {1,2}
    \draw (-0.05,\y) -- (0.05,\y)
      node[left] {$\y$};

\end{tikzpicture}
\end{center}

\section* {第十一题}

$$y_1[n] = x[2n], y_2[n] = \begin{cases}
  x[\frac{n}{2}], & n \equiv 0 (mod 2), \\
  0, & n \equiv 1 (mod 1).
\end{cases}$$

\subsection{$x[n]$周期$T_0 \longrightarrow y_1[n] \quad ?$}

正确，其基波周期为

$$ T_1 = \begin{cases}
  T_0, &T_0 \equiv 1(mod 2), \\
  \frac{T_0}{2}, &T_0 \equiv 0(mod 2).
\end{cases}$$

\subsection{$y_1[n]$周期$T_1 \longrightarrow x[n] \quad ?$}

错误。不妨取$y_1[n] \equiv 0$为周期信号。定义
$$ x[n] = \begin{cases}
  0, &n\equiv 0(mod 2), \\
  n, &n\equiv 1(mod 2).
\end{cases}$$
显然不为周期的，因为$y[n]$不稳定。

\subsection{$x[n]$周期$T_0 \longrightarrow y_2[n] \quad ?$}

正确，其基波周期为
$$T_2 = \begin{cases}
  1 & x[n] \equiv 0, \\
  max(2, 2T_0) & x[n] \ne 0,
\end{cases}$$

\subsection{$y_2[n]$周期$T_2$ $\longrightarrow x[n] \quad ?$}

正确，其基波周期为
$$T_0 = \begin{cases}
  1 & y_2[n] \equiv 0, \\
  \frac{T_2}{2} & y_2[n] \equiv 0.
\end{cases}$$

\section* {第十二题}

$x(t)=e^{j\omega_0 t}$,基波频率$\omega_0$,基波周期$T_0=\frac{2\pi}{\omega_0}$,等距采样得到$x[n]=e^{j\omega_0nT}$

\subsection{$x[n]$是周期的$\Leftrightarrow \frac{T}{T_0} \in \mathbb{Q}$}

$\qquad x[n]=e^{j\omega_0nT} 
\quad$是周期的$\quad \Leftrightarrow \quad\frac{\omega}{2\pi} = \frac{\omega_0 T}{2\pi} = \frac{T}{T_0}\in \mathbb{Q}$

\subsection{$x[n]$是周期的，即$\frac{T}{T_0}=\frac{p}{q}$时，求基波频率及基波周期。基波频率表示成$\omega_0T$的分式}

假设$(p, q)$既约。

$$T_0 = \frac{2\pi}{\omega_0} \Rightarrow T = \frac{p}{q} T_0 = \frac{2p\pi}{q\omega_0} \Rightarrow \frac{\omega}{2\pi} = \frac{\omega_0 T}{2\pi} = \frac{p}{q}$$

因此，基波周期$T'=q$。基波频率$\omega'=\frac{2\pi}{T'}=\frac{2\pi}{q}=\frac{\omega_0 T}{p}$ 

\subsection{同上假设。确定需要多少个$x(t)$的周期才能得到$x[n]$的一个样本}

$$\frac{\omega_0 T T'}{2\pi} = \frac{\omega_0 \frac{pT_0}{q} q}{2\pi} = p$$

\section* {第十三题}

对系统$S$，输入$x[n]$与输出$y[n]$关系为
$$y[n] = x[n]\{g[n]+g[n-1]\}$$

\subsection{$g[n] \equiv 1$, 证明$S$是时不变的}

$$y[n] = 2x[n]$$

$\forall n_0 \in \mathbb{R}, x_2[n] = x_1[n-n_0]$， 有：
$$y_2[n] = 2x_2[n] = 2x_1[n-n_0] = y_1[n-n_0]$$

\subsection{$g[n] = n$,证明$S$不是时不变的}

$$y[n] = (2n-1)x[n]$$

对于$x_2[n] = x_1[n-1]$，有：
$$y_2[n] = (2n-1)x_2[n] = (2n-1)x_1[n-1]$$ 
$$y_1[n-1] = (2n-3)x[n-1] \ne y_2[n]$$

\subsection{$g[n] = 1+(-1)^n$，证明$S$是时不变的}

$$y[n] = x[n]\{(1+(-1)^n) + (1+(-1)^{n+1})\} = 2x[n]$$

与第一问的形式一致，结论一致

\section* {第十四题}

以下仅证明连续系统的情况。离散系统同理。

\subsection{两个线性时不变系统的级联仍线性时不变？}

正确。记两个系统分别为：

$$S_1: \quad y(t) = f(x(t))$$
$$S_2: \quad y(t) = g(x(t))$$

不妨设$S_1$输出为$S_2$输入，则其级联为 $y(t) = g(f(x(t)))$。 \\

首先验证线性性。

$$\forall a, b \in \mathbb{R}$$

\begin{align}
  y(t) &= gf(x(t)) = gf(ax_1(t)+bx_2(t)) \\
       &= g(af(x_1(t)) + bf(x_2(t))) \\
       &= agf(x_1(t)) + bgf(x_2(t)) \\
       &= ay_1(t) + by_2(t)
\end{align}

然后验证时不变性。

$$\forall t_0 \in \mathbb{R}, x_2(t) = x_1(t-t_0)$$

记中间信号及最终输出为
\[
v_i(t)=S_1[x_i](t),\qquad
y_i(t)=S_2[v_i](t),\qquad i=1,2.
\]

由 $S_1$ 的时不变性，
\[
v_2(t)=S_1[x_1(\cdot-t_0)](t)
      =S_1[x_1](t-t_0)
      =v_1(t-t_0).
\]

再由 $S_2$ 的时不变性，
\[
y_2(t)=S_2[v_1(\cdot-t_0)](t)
      =S_2[v_1](t-t_0)
      =y_1(t-t_0).
\]



\subsection{两个非线性系统的级联仍非线性？}

错误。构造以下两个非线性系统：

$$S_1: \quad y(t) = \sqrt[3]{x(t)}$$
$$S_2: \quad y(t) = (x(t))^3$$

其串联显然线性。

\end{document}